%% kapa: exam=1 round=30 year=2019 period=1 form=A subject=real_estate subject_ko=부동산학원론 seq=39 number=119
%% source: Q-Net 감정평가사 기출문제 (cst00309) artlSeq=5205205 file=extracted/1008/5205205/2019년 제30회 감정평가사 1차시험 문제지/2019년 제30회 감정평가사 1차시험 문제지/2019년도 제30회 감정평가사 1차 1교시 A형.hwp
%% conversion: kordoc 4.24.51 -> tools/gen1.py (auto); review=proofread-sonnet; flags=none
\begin{question}{119}
지대이론 및 도시공간구조이론에 관한 설명으로 옳지 않은 것은?
\begin{choices}
  \item 리카도(D. Ricardo)는 비옥한 토지의 희소성과 수확체감의 법칙으로 인해 지대가 발생한다고 보았다.
  \item 마샬(A. Marshall)은 일시적으로 토지와 유사한 성격을 가지는 생산요소에 귀속되는 소득을 준지대로 보았다.
  \item 알론소(W. Alonso)는 각 토지의 이용은 최고의 지대지불의사가 있는 용도에 할당된다고 보았다.
  \item 호이트(H. Hoyt)는 저급주택지가 고용기회가 많은 도심지역과의 교통이 편리한 지역에 선형으로 입지한다고 보았다.
  \item 해리스(C. Harris)와 울만(E. Ullman)은 도시 내부의 토지이용이 단일한 중심이 아니라 여러 개의 전문화된 중심으로 이루어진다고 보았다.
\end{choices}
\end{question}
\answer{4}
